\documentclass[12pt,a4paper]{report}

\usepackage[spanish]{babel}
\usepackage[utf8]{inputenc}
\usepackage[T1]{fontenc}

\begin{document}


\chapter{Introducción}

La gestión y evaluación de la calidad de las aguas subterráneas constituye un desafío estratégico para el desarrollo sostenible, la planificación territorial y la protección ambiental. En este contexto, el uso de tecnologías de sistemas de información geográfica (SIG) y plataformas web de mapas son herramientas útiles para la visualización, el análisis y la difusión de información geoespacial. El presente trabajo de tesis se orienta al diseño y desarrollo de una plataforma web interactiva de mapas aplicada al análisis de la calidad del agua subterránea del acuífero Patiño, con énfasis en su reutilización en otros escenarios similares.

\section{Antecedentes y problemática}

Las plataformas web de mapas permiten la visualización, exploración y manipulación de datos geoespaciales mediante interfaces accesibles a través de navegadores web, sin requerir software especializado (REF). A nivel internacional, existen múltiples experiencias orientadas al monitoreo y gestión de aguas subterráneas que integran redes piezométricas, puntos de muestreo y datos de calidad de agua, utilizando tanto tecnologías propietarias como soluciones de código abierto. Estas plataformas han demostrado ser herramientas eficaces para apoyar la toma de decisiones en gestión ambiental, planificación territorial y monitoreo de recursos hídricos.

En Paraguay, si bien se han desarrollado diversos visores de mapas web en instituciones públicas y académicas con fines informativos y de planificación, la mayoría de estas soluciones presentan limitaciones significativas. Entre ellas se destacan la dependencia de software propietario, los altos costos de licenciamiento, la limitada sostenibilidad tecnológica y la escasa interactividad para el análisis avanzado de datos. Asimismo, muchos visores se restringen a la visualización estática de información, sin permitir consultas dinámicas, integración de datos actualizados o participación activa del usuario.

En la Facultad Politécnica de la Universidad Nacional de Asunción (FPUNA) se implementó en años anteriores un visor de mapas web para la socialización de resultados de investigaciones vinculadas al acuífero Patiño. No obstante, dicha solución fue desarrollada con herramientas de código abierto que posteriormente quedaron descontinuadas, lo que imposibilita su mantenimiento, actualización y reutilización.

Ante este escenario, surge la necesidad de desarrollar una nueva plataforma web de mapas que utilice tecnologías de los Sistemas de Información Geográfica (SIG) actuales, abiertas y sostenibles, que permita no solo la visualización, sino también la gestión, análisis y cálculo automático de indicadores de calidad de agua subterránea.

\section{Justificación y motivación}

La importancia del acuífero Patiño como fuente de abastecimiento de agua para Asunción y el área metropolitana hace imprescindible contar con herramientas modernas que faciliten el acceso, análisis y comprensión de la información relacionada con su calidad. Una plataforma web interactiva de mapas, basada en software de código abierto, permitiría reducir costos de implementación, evitar la dependencia de proveedores propietarios y promover la interoperabilidad mediante estándares abiertos.

El presente trabajo contribuye al fortalecimiento de capacidades locales en el desarrollo de soluciones Web-GIS, fomentando la reutilización y adaptación de la plataforma en otros contextos, como el monitoreo de otros acuíferos, aguas superficiales, planificación territorial o gestión de riesgos ambientales. Asimismo, la motivación del trabajo se sustenta en la necesidad de superar las limitaciones de soluciones previas, garantizando la sostenibilidad tecnológica mediante la adopción de herramientas con comunidades activas de desarrollo.

\section{Objetivo general}

Desarrollar e implementar una plataforma web de mapas, basada en herramientas de código abierto, que posibilite la geolocalización en la evaluación de la calidad de aguas subterráneas, permitiendo la digitalización, almacenamiento, visualización y análisis de datos geoespaciales para mejorar la toma de decisiones y su reutilización en otros escenarios similares.

\section{Objetivos específicos}

\begin{itemize}
	\item Analizar el estado del arte de las tecnologías de sistemas de información geográfica de código abierto existentes para la publicación de mapas web.
	\item Desarrollar una plataforma web interactiva para la geolocalización, visualización y carga dinámica de mapas.
	\item Establecer diferentes categorías de usuarios para la utilización de la plataforma.
	\item Proponer una solución de código abierto, escalable y reutilizable en otros escenarios similares.
	\item Recolectar, filtrar y normalizar los datos necesarios para la creación de una base de datos geoespacial de muestreos de calidad de agua.
	\item Desarrollar un módulo que permita el cálculo automático de la calidad del agua subterránea a partir de parámetros fisicoquímicos y microbiológicos.
\end{itemize}

\section{Limitaciones}

El alcance del presente trabajo se limita al desarrollo y validación de una plataforma web de mapas aplicada al acuífero Patiño, utilizando conjuntos de datos disponibles y seleccionados para fines académicos. No se contempla la realización de campañas de muestreo de campo, sino la utilización de datos existentes. Asimismo, los resultados obtenidos dependerán de la calidad, disponibilidad y actualización de la información utilizada como entrada en la plataforma.

\section{Organización del trabajo}

El documento de tesis se estructura en varios capítulos. El Capítulo 1 presenta la introducción, los antecedentes, la justificación, los objetivos, las limitaciones y la organización general del trabajo. El Capítulo 2 aborda el marco teórico y el estado del arte relacionado con las plataformas Web-GIS y el uso de tecnologías de código abierto. El Capítulo 3 describe la metodología empleada para el desarrollo de la plataforma. El Capítulo 4 presenta el diseño e implementación de la solución propuesta. El Capítulo 5 expone los resultados obtenidos y su análisis. Finalmente, el Capítulo 6 presenta las conclusiones y recomendaciones para trabajos futuros.

\end{document}

